\subsection{A source quotient detected by the CM depth-two measure}
\label{subsec:cm-measure-class-two}

Assume \(\mathcal X_2^+\) is finite and retain
\(H=\operatorname{Gal}(\Omega/L_\infty)\) and
\(A=H^{\mathrm{ab}}\).  The proved CM results make \(H\)
free pro-\(p\) and make the joint elliptic Soul\'e map
injective on \(A\).

There is a second useful measure realization of this map.
Put \(V=T_p(E)\simeq\mathbb Z_p^2\), written additively.
By \cite[Lemma 2.4, Definition 2.5 and Theorem 2.6,
pp.~7--8]{Ishii}, the canonical boundary-preserving
lift of the outer action gives
\[
 \alpha:A\longrightarrow
       \mathbb Z_p[[V]](1).
\]
Here \(\mathbb Z_p[[V]]\) denotes bounded measures on \(V\).
If \(U_i=\log(1+T_i)\), its exponential transform is
\[
 f_h(U_1,U_2)
 =\sum_{\substack{r,s\geq0\\r+s\geq2\ {\rm even}}}
 \frac{\kappa_{(r+1,s+1)}(h)}
      {(1-p^{r+s})r!s!}\,U_1^rU_2^s .
\]
The coefficients outside the admissible CM weight set
vanish.  The already proved joint-kernel equality makes
\(\alpha\) injective.

Define two automorphisms of \(V\times V\) by
\[
 t_1(x,y)=(x,x+y),\qquad t_2(x,y)=(x+y,y),
\]
and define the bounded-measure operator
\[
 \mathcal J=(t_1)_*+(t_2)_*-\operatorname{id}.
\]
The one-dimensional restriction of this operator occurs in
\cite[equations (29) and (32), p.~233]{LuqueNovelliThibon}.
To specify the comparison, send their coefficient
\(\phi_{n+1}\) to \(U^n/n!\), so that
\(\Phi(t)\) becomes \(\exp(tU)\), the exponential transform
of the point measure \(\delta_t\). Send an ordinary bracket
\([f,g]\) to \(f(U)g(V)-g(U)f(V)\). Their equation (32)
then becomes
\[
 \delta_s\wedge\delta_{s+t}
 +\delta_{s+t}\wedge\delta_t-\delta_s\wedge\delta_t
                 =\mathcal J(\delta_s\wedge\delta_t).
\]
For every \(v\in V\), the map \(i_v(t)=tv\) preserves
addition and gives
\((i_v\times i_v)_*\mathcal J_1
 =\mathcal J(i_v\times i_v)_*\).
The source calculation below proves the displacement formula
for the CM arithmetic action on the two-dimensional Tate module.

Use the alternating tensor inclusion to define
\[
 J_A:\widehat{\bigwedge}^2 A
 \longrightarrow\mathbb Z_p[[V\times V]](2),\qquad
 a\wedge b\longmapsto
 \mathcal J\bigl(\alpha(a)\otimes\alpha(b)
                         -\alpha(b)\otimes\alpha(a)\bigr).
\]
Let \(N_2^{\rm CM}\) be the inverse image of
\(\ker J_A\) under the canonical map
\(\gamma_2H\to\gamma_2H/\gamma_3H
 \simeq\widehat{\bigwedge}^2A\).

\begin{proposition}\label{prop:cm-decomposable-commutators}
For \(a,b\in A\) whose images in \(A[1/p]\) are linearly
independent, \(J_A(a\wedge b)\ne0\). Thus every such pair
has a nonzero commutator in \(Q_2^{\rm CM}\).
\end{proposition}
\begin{proof}
Write \(f=\alpha(a)\), \(g=\alpha(b)\) in the exponential
coordinates above. Both have zero constant and linear terms.
The exponential transform is injective: the Amice transform
identifies bounded measures with integral power series, and
the substitution \(T_i=\exp(U_i)-1\) has the inverse
\(U_i=\log(1+T_i)\) over \(\mathbb Q_p\).
The transform of \(J_A(a\wedge b)\) is
\begin{align*}
 &f(u)g(u+v)-g(u)f(u+v)\\
 &\quad+f(u+v)g(v)-g(u+v)f(v)-f(u)g(v)+g(u)f(v).
\end{align*}
Its coefficient linear in \(v_i\) is
\(f(u)\partial_i g(u)-g(u)\partial_i f(u)\), since every
term containing \(f(v)\) or \(g(v)\) has order at least two.
If the displayed transform vanishes, both derivatives of
\(g/f\) vanish in \(\operatorname{Frac}(\mathbb Q_p[[U_1,U_2]])\).
Embed this field into \(\mathbb Q_p((U_1))((U_2))\).
Every nonzero power series is invertible there after extracting
its first nonzero \(U_2\)-coefficient; the embedding commutes
with formal derivatives. The \(U_2\)-derivative forces the
ratio to belong to \(\mathbb Q_p((U_1))\), and its
\(U_1\)-derivative forces it to belong to \(\mathbb Q_p\),
by coefficient comparison in characteristic zero. Thus
\(g=cf\). Injectivity of \(\alpha[1/p]\) gives \(b=ca\),
contradicting independence. The commutator assertion follows
from the quotient theorem below.
\end{proof}

\begin{theorem}\label{thm:cm-measure-class-two}
The subgroup \(N_2^{\rm CM}\) is closed and stable under
the arithmetic action, and contains the outer
kernel.  The quotient
\[
                       Q_2^{\rm CM}=H/N_2^{\rm CM}
\]
is torsion-free of class at most two, has abelianization
\(A\), and has central commutator module
\(\operatorname{im}J_A\).
Its class is exactly two for every class-number-one
CM field in the stated range.
\end{theorem}

\begin{proof}
Let \(\Pi\) be the geometric pro-\(p\) fundamental group
of \(E-\{O\}\), with the generators \(x_1,x_2\)
of \cite[Lemma 2.1]{Ishii}. Put \(C=[\Pi,\Pi]\), \(C_r=\gamma_rC\), and
\(M=C/C_2=\mathbb Z_p[[V]]\bar z\), with
\(z=[x_2,x_1]\) as in \cite[Lemma 2.1]{Ishii}.
The closed subgroup \(C\) of the free pro-\(p\) group
\(\Pi\) is free pro-\(p\).  The integral class-two Hall
basis, first on finite subsets of a convergent free
basis and then by inverse limits, gives
\[
 C_2/C_3=\widehat{\bigwedge}^{\,2}M,
 \qquad \gamma_2H/\gamma_3H=
            \widehat{\bigwedge}^{\,2}A.
\]
The exterior square identifies with the antisymmetric
submodule of the completed tensor square by
\(a\wedge b\mapsto a\otimes b-b\otimes a\).
Since \(p\) is odd, half the exterior projection is
its inverse on this submodule.  Thus
\(C_2/C_3\) is a closed submodule of
\(\mathbb Z_p[[V\times V]](2)\).

For \(v\in V\), write \(e_v=v\bar z\), and let
\(T_v\) denote translation on \(M\) and simultaneous
translation on its tensor square.  Let \(f=s(h)\)
be the canonical special lift of \(h\in H\), and set
\(a=\alpha(h)\bar z\).  This lift fixes \(z\) exactly:
its exponent on the inertia generator is its
abelian determinant, which is one.  For \(q\in\Pi\)
with abelian image \(v\), put
\(k_f(q)=f(q)q^{-1}\in C\).
The crossed-homomorphism identity and Ishii's defining
formula on \(x_1,x_2\) give
\[
                 k_f(q)\bmod C_2=(T_v-1)a.
\]
This follows first on words and then by continuity.
Because \(C\) is normally generated by \(z\), the identity
\[
 f(qzq^{-1})=k_f(q)(qzq^{-1})k_f(q)^{-1}
\]
shows that \(f\) acts trivially on \(M\).  Its first
nontrivial action on \(C\) is the continuous map
\[
 \delta_a:M\longrightarrow C_2/C_3,
 \qquad \delta_a(e_v)=[T_va-a,e_v].
\]
In particular, it depends only on \(\alpha(h)\), and
\begin{equation}\label{eq:cm-delta-translation}
 \delta_a(T_vb)=T_v\delta_a(b)+[T_va-a,T_vb].
\end{equation}
The formula follows on point measures by substitution
and extends to all bounded measures by continuity.

Take another lift \(g=s(k)\), with
\(b=\alpha(k)\bar z\).  Their metabelian actions
commute.  In the class-two group \(C/C_3\), the BCH
law involves only the invertible denominator \(2\).
Use
\(k_{fg}(q)=f(k_g(q))k_f(q)\) and its reversed-order
counterpart.  Their unknown central displacements
cancel, and the surviving terms are
\begin{align*}
 k_{[f,g]}(q)\bmod C_3
  ={}&\delta_a((T_v-1)b)-\delta_b((T_v-1)a)\\
     &+[(T_v-1)b,(T_v-1)a].
\end{align*}
Indeed a special automorphism with all displacements
in \(C_2\) acts trivially on \(C/C_3\), by the
conjugate-of-\(z\) identity.  Hence replacing \(q\)
by \(gf(q)\) when passing from \(fg,gf\) to
\([f,g]\) does not change this displayed displacement.
Writing \(A_0=T_va,B_0=T_vb\), the bracket terms
introduced by \eqref{eq:cm-delta-translation} simplify as
\[
 [A_0-a,B_0]-[B_0-b,A_0]+[B_0-b,A_0-a]
                         =[A_0,B_0]-[a,b].
\]
We therefore obtain the exact identity
\[
 k_{[f,g]}(q)\bmod C_3
       =(T_v-1)(\delta_ab-\delta_ba+[a,b]).
\]
From the point-measure formula for \(\delta\),
\begin{align*}
 \delta_ab-\delta_ba+[a,b]
 &=\iint\bigl([e_x,e_{x+y}]+[e_{x+y},e_y]
                         -[e_x,e_y]\bigr)\,da(x)\,db(y)\\
 &=\mathcal J(a\wedge b).
\end{align*}
This proves the integral group comparison
\begin{equation}\label{eq:cm-integral-J-comparison}
 k_{[s(h),s(k)]}(q)\bmod C_3
      =(T_v-1)J_A(h^{\rm ab}\wedge k^{\rm ab}).
\end{equation}

For completeness, let \(\mathcal D_r\) be the special
automorphisms acting trivially on \(V\) whose
displacements all lie in \(C_r\).
An element of \(\mathcal D_1\) acts trivially on
\(C_2/C_3\), whereas an element of \(\mathcal D_2\)
acts trivially on \(C/C_3\).  Comparing their product
displacements and using \([C,C_2]\subset C_3\) proves
\([\mathcal D_1,\mathcal D_2]\subset\mathcal D_3\).
Consequently \(s(\gamma_3H)\subset\mathcal D_3\),
and additivity of displacement on \(\mathcal D_2\)
extends \eqref{eq:cm-integral-J-comparison} continuously
to every \(h\in\gamma_2H\):
\[
 k_{s(h)}(q)\bmod C_3=(T_v-1)J_A(h\bmod\gamma_3H).
\]
For \(q=x_1\), the operator \(T_v-1\) is multiplication
by
\[
 (1+T_1)(1+S_1)-1=T_1+S_1+T_1S_1
\]
in the power-series domain
\(\mathbb Z_p[[T_1,T_2,S_1,S_2]]\), so it is injective.
An outer-kernel element belongs to \(\gamma_2H\) by
abelian detection and has \(s(h)=1\), since the
canonical special lift is inverse to the outer
projection in \(F^3\).  The last identity therefore
places it in \(N_2^{\rm CM}\).

Here is also the exact arithmetic compatibility.
An arithmetic outer automorphism has a representative
preserving \(\langle z\rangle\); any two such
representatives differ by a power of
\(\operatorname{inn}(z)\), which commutes with every
\(s(h)\).  Uniqueness in \(F^3\) makes the special lift
equivariant.  On \(M\), conjugation is the linear
Tate action with determinant twist.  Both shears
commute with simultaneous linear changes on \(V\),
so \(J_A\) is equivariant with the displayed second
twist.  Thus its inverse-image kernel is stable
under the arithmetic action.  It is closed, and its
image is compact.  The resulting exact central extension
\[
 1\longrightarrow\operatorname{im}J_A
 \longrightarrow Q_2^{\rm CM}\longrightarrow A
 \longrightarrow1
\]
has commutator subgroup exactly \(\operatorname{im}J_A\).
Both end modules are \(p\)-torsion-free, so the group
is torsion-free and has the claimed abelianization.

It remains to prove that the commutator module is nonzero.
Let \(I_K\) denote the admissible weights, and put
\[
 \mathcal S=\{\boldsymbol m\in I_K:
  1\leq m_i\leq4,\ (5,5)-\boldsymbol m\in I_K\}.
\]
These are exactly the moments that can contribute
to the bidegree-\((5,5)\) bracket coefficient.
The ray-class description and two CM Tate characters
in \cite[pp.~4--5]{Ishii} show that
\(\Gamma=\operatorname{Gal}(L_\infty/L_1)\) has open
rank-two principal-unit image. Write
\(\chi_{\boldsymbol m}=\chi_1^{m_1}\chi_2^{m_2}\)
for the corresponding product of the two Tate characters. Choose \(\gamma\)
so that \(\chi_{\boldsymbol m}(\gamma)\), for
\(\boldsymbol m\in\mathcal S\), are pairwise distinct.
This is possible because after the \(p\)-adic logarithm
the finitely many equality loci are proper linear
subspaces of an open \(\mathbb Z_p^2\).
For \(\boldsymbol m_0=(2,2),(3,3)\), define the integral
finite-difference operator
\[
 P_{\boldsymbol m_0}(X)=
 \prod_{\substack{\boldsymbol m\in\mathcal S\\
                   \boldsymbol m\ne\boldsymbol m_0}}
       (X-\chi_{\boldsymbol m}(\gamma)).
\]
Both \(\kappa_{22}\) and \(\kappa_{33}\) are nonzero
by \cite[Theorem 2.8(2) and the following Soul\'e-finiteness
paragraph, p.~9]{Ishii}.
Choose \(a_0,b_0\in A\) on which these respective
characters are nonzero, and set
\(a=P_{22}(\gamma)a_0,b=P_{33}(\gamma)b_0\).
Equivariance gives
\[
 \kappa_{\boldsymbol m}(P(\gamma)d)
   =P(\chi_{\boldsymbol m}(\gamma))\kappa_{\boldsymbol m}(d).
\]
Thus all unwanted moments in \(\mathcal S\) vanish,
while the two desired moments remain nonzero.

In exponential moment variables \(U=(x,y),V=(z,t)\),
\(\mathcal JF=F(U+V,V)+F(U,U+V)-F(U,V)\).
Direct multiplication gives
\[
 J(XY,X^2Y^2)
       =(xy-zt)(xt+zt+yz)(xt+xy+yz)\ne0.
\]
The bidegree-\((3,3)\) polynomial component of
\(J_A(a\wedge b)\), corresponding to bracket weight
\((5,5)\), is this polynomial multiplied by
\[
 \frac{\kappa_{22}(a)\kappa_{33}(b)}
      {(1-p^2)(1-p^4)(2!)^2}\ne0.
\]
The definition of \(\mathcal S\) excludes every other
contribution.  Hence \(\operatorname{im}J_A\ne0\),
proving exact class two for every \(K\) in the statement.
\end{proof}

\subsection{Centralizers and free class-two source subgroups}
\label{subsec:cm-class-two-subgroups}

Retain the hypotheses and notation of
Theorem~\ref{thm:cm-measure-class-two}, and write
\[
 Q=Q_2^{\rm CM},\qquad Z=\operatorname{im}J_A,
 \qquad \pi:Q\longrightarrow A.
\]
The following statements use the proved integral map from the
actual source group to this quotient.

\begin{proposition}\label{prop:cm-class-two-centralizers}
The centre of \(Q\) is exactly \(Z\).  For every \(q\in Q\)
with \(a=\pi(q)\ne0\),
\[
 C_Q(q)=\pi^{-1}\bigl(A\cap\mathbb Q_p a\bigr).
\]
The intersection is taken inside \(A[1/p]\), is a free
\(\mathbb Z_p\)-module of rank one, and the displayed
centralizer is, noncanonically, \(Z\times\mathbb Z_p\).
In particular, every abelian subgroup of \(Q\) has image
of \(\mathbb Z_p\)-rank at most one in \(A\).
\end{proposition}

\begin{proof}
The commutator of two elements with abelian images \(a,b\)
is \(J_A(a\wedge b)\).  Proposition~\ref{prop:cm-decomposable-commutators}, together
with injectivity of \(\alpha[1/p]\), gives the exact criterion
\[
 J_A(a\wedge b)=0\quad\Longleftrightarrow\quad
 b\in\mathbb Q_pa\qquad(a\ne0).
\]
This proves the asserted centralizer formula.

Here is the integral topology of the line intersection.
Choose a continuous \(\mathbb Z_p\)-linear functional
\(\ell:A\to\mathbb Z_p\) with \(\ell(a)\ne0\).
Such a functional exists because the compact torsion-free
abelian pro-\(p\) group \(A\) is a product of copies of
\(\mathbb Z_p\).  If \(\lambda a\in A\), then
\(\lambda\ell(a)\in\mathbb Z_p\), which bounds the
denominator of \(\lambda\).  The allowed scalars form a
nonzero closed \(\mathbb Z_p\)-submodule of this bounded
rank-one lattice.  They are consequently \(p^e\mathbb Z_p\)
for some integer \(e\).  Thus \(A\cap\mathbb Q_pa\)
has a generator \(a_0\).  Choose \(q_0\in Q\) above
\(a_0\).  The continuous powers of \(q_0\) identify
its closed procyclic group with \(\mathbb Z_p\), and
\((z,t)\mapsto zq_0^t\) is an isomorphism from
\(Z\times\mathbb Z_p\) onto the displayed centralizer.

The module \(A[1/p]\) has dimension at least two: for
example, the moment-isolated elements constructed below
have distinct nonzero initial homogeneous terms.
Every nonzero \(a\) therefore has a \(\mathbb Q_p\)-independent
\(b\), so an element above \(a\) is not central.
Hence the centre is precisely \(Z\).  The assertion about
abelian subgroups follows from the same centralizer formula.
\end{proof}

\begin{corollary}\label{cor:cm-three-generator-class-two}
Let \(r\le3\), and let \(h_1,\ldots,h_r\in H\) have
\(\mathbb Q_p\)-linearly independent images in \(A[1/p]\).
For the closed subgroup \(U\) they generate, the map
\(U\to Q\) has kernel exactly \(\gamma_3U\), and
\(U/\gamma_3U\) is the free class-two pro-\(p\) group
on these \(r\) generators.
\end{corollary}

\begin{proof}
Let \(B\) be the \(\mathbb Z_p\)-span of their abelian
images.  Every bivector in the exterior square of a
vector space of dimension at most three over
\(\mathbb Q_p\) is decomposable.  The preceding
derivative argument thus makes
\(J_A:\bigwedge^2B\to Z\) injective.
Map the free pro-\(p\) group \(P_r\) to \(H\) using
the chosen elements.  The composite to \(Q\) kills
\(\gamma_3P_r\); its maps on the abelian quotient and
on \(\gamma_2P_r/\gamma_3P_r\simeq\bigwedge^2\mathbb Z_p^r\)
are injective.  Its kernel is therefore exactly
\(\gamma_3P_r\).  The image of that subgroup in \(U\)
is \(\gamma_3U\), which proves both conclusions.
\end{proof}

\begin{proposition}\label{prop:cm-countable-class-two}
There is a closed subgroup \(U\subset H\), generated
by a sequence tending to one, for which
\[
 U\cap N_2^{\rm CM}=\gamma_3U,
 \qquad U/\gamma_3U\simeq P/\gamma_3P,
\]
where \(P\) is the free pro-\(p\) group on a countably
infinite convergent sequence of generators.
Thus the complete free class-two quotient of this
actual source subgroup is detected.  Every finite
initial subsequence gives the corresponding conclusion
in finite rank.
\end{proposition}

\begin{proof}
For \(i\ge1\), put \(k_i=2^i\) and
\(d_i=2k_i-2\).  The character
\(\kappa_{(k_i,k_i)}\) is nonzero by
\cite[Theorem 2.8(2) and the Soul\'e-finiteness paragraph,
p.~9]{Ishii}; this is the ordinary \(k_i\)-th
Tate twist, with \(k_i\ge2\).

Let
\[
 \mathcal S_i=\{\boldsymbol m\in I:
                              |\boldsymbol m|\le2k_i\}.
\]
Put \(\Gamma_2=\operatorname{Gal}(L_\infty/L_1)\), and choose \(\gamma_i\in\Gamma_2\) for which the finitely
many characters \(\chi_{\boldsymbol m}(\gamma_i)\),
\(\boldsymbol m\in\mathcal S_i\), are pairwise distinct.
The open rank-two principal-unit image of the two CM
Tate characters, established by the ray-class morphism
in \cite[pp.~4--5]{Ishii}, permits this choice:
the forbidden loci after taking logarithms are a finite
union of proper linear subspaces.  Put
\[
 P_i(X)=\prod_{\substack{\boldsymbol m\in\mathcal S_i\\
                         \boldsymbol m\ne(k_i,k_i)}}
                 (X-\chi_{\boldsymbol m}(\gamma_i)).
\]
Choose \(b_i^0\in A\) with
\(\kappa_{(k_i,k_i)}(b_i^0)\ne0\), and set
\(b_i=P_i(\gamma_i)b_i^0\).  Equivariance shows that
all its moments in \(\mathcal S_i\), except the desired
one, vanish, and that the desired one remains nonzero.

Choose a lift \(\widetilde h_i\in H\) of \(b_i\), and
put \(h_i=\widetilde h_i^{p^i}\), with abelian image
\(a_i=p^ib_i\).  This sequence tends to one: in every
finite quotient of \(H\), a sufficiently large power
\(p^i\) kills every element.  The exponential transform
\(f_i\) of \(\alpha(a_i)\) has the form
\[
 f_i(U_1,U_2)=c_i(U_1U_2)^{k_i-1}
                    +\text{terms of total degree greater than }d_i,
 \qquad c_i\in\mathbb Q_p^\times.
\]
In particular, the \(d_i\) are strictly increasing.

Let \(P\) be free on the convergent sequence
\(x_i\to1\).  Its abelianization is
\(\prod_{i\ge1}\mathbb Z_p x_i\), and the integral
class-two Hall calculation gives
\[
 \gamma_2P/\gamma_3P
       =\prod_{i<j}\mathbb Z_p[x_i,x_j].
\]
These assertions follow by taking inverse limits of
the corresponding finite-generator, finite-quotient
Hall presentations.  The continuous map \(x_i\mapsto h_i\)
induces an injective abelian map: in a nonzero completed
linear combination of the \(f_i\), the least index
with nonzero coefficient supplies a unique nonzero
initial degree.

The initial homogeneous term of \(J(f_i,f_j)\) is
\[
 c_ic_j J\bigl((U_1U_2)^{k_i-1},
              (U_1U_2)^{k_j-1}\bigr),
\]
of degree
\[
 d_i+d_j=2^{i+1}+2^{j+1}-4.
\]
It is nonzero by the derivative computation above:
the two displayed monomials have nonconstant ratio.
These degrees are pairwise distinct for \(i<j\), by
uniqueness of binary expansion.  A nonzero completed
alternating combination therefore again has a unique
least initial degree, which cannot cancel.
For each total degree there are only finitely many
pairs to consider, so this argument applies to the
completed product, not only to finite combinations.
Consequently the composite
\(P\to H\to Q\) induces an injection on both its
abelian quotient and its class-two central quotient.
Its kernel is exactly \(\gamma_3P\).

Let \(U\) be the image of \(P\) in \(H\).  It is compact
and hence closed.  The image of \(\gamma_3P\) is
\(\gamma_3U\), by the commutator definition and
compactness.  The exact kernel just calculated gives
\(U\cap N_2^{\rm CM}=\gamma_3U\) and the asserted
isomorphism.  No injectivity assertion for the full
map \(P\to H\) is needed for these conclusions.
\end{proof}

The subgroups in Proposition~\ref{prop:cm-countable-class-two}
depend on the moment-isolating choices.  The surrounding
quotient \(Q_2^{\rm CM}\) remains the canonical arithmetic
quotient of Theorem~\ref{thm:cm-measure-class-two}.

